\documentclass[10pt,a4paper]{amsart}
\usepackage[margin=19mm]{geometry}
\usepackage{amsmath,amssymb,mathtools}
\usepackage[scr=rsfs,cal=euler]{mathalfa}
\usepackage{xfrac}
\usepackage[unicode,hidelinks,bookmarks=false]{hyperref}
\newcommand{\ie}{i.e.\ }
\newcommand{\cf}{cf.\ }
\newcommand{\resp}{respectively\ }
\newcommand{\Subsection}{\subsection}
\providecommand{\nobreakdash}{\nobreak\hbox{-}}
\setlength{\emergencystretch}{2em}
\multlinegap0pt
\usepackage{amssymb}
\usepackage{tikz}
\usepackage{xfrac}
\usepackage[shortlabels]{enumitem}
\usepackage{mathtools}
\usepackage{float}
\newtheorem{lemma}{Lemma}
\newcommand{\defeq}{\vcentcolon=}
\renewcommand{\geq}{\geqslant}
\renewcommand{\leq}{\leqslant}
\renewcommand{\wp}{\ensuremath\raisebox{\depth}{$\mathchar"17D$}}
\title{Étale extensions of unipotent torsors}
\author{Gabriel Bassan}
\date{Source: arXiv:2605.04182v1 (CC BY 4.0)}
\begin{document}
\maketitle

\noindent\emph{Source and licence.} Étale extensions of unipotent torsors. Version arXiv:2605.04182v1 (CC BY 4.0), distributed by arXiv under the Creative Commons Attribution 4.0 International licence, http://creativecommons.org/licenses/by/4.0/, as recorded in the arXiv licence metadata for this exact version and shown on the abstract page https://arxiv.org/abs/2605.04182v1. Full licence notice: \texttt{licenses/arxiv-2605.04182-CC-BY-4.0.txt}.\par\smallskip

\noindent\textit{Editorial scope.} Counted unit: the article's lemma pi (source0.tex lines 554--560). The article's complete original proof of it is delivered with it (source0.tex lines 561--606, one \texttt{proof} environment of 46 lines). No mathematics is written, repaired or re-derived: the statement and the proof are the article's own lines, in the article's own notation, and no retained line is substituted or re-flowed.  Retained uncounted as the source-local context the proof needs: lines 540--543.  Nothing was omitted from any retained span, so the package carries no omission record.  No retained line contains a citation, so the wrapper carries no bibliography and no bibliography entry.  No retained line contains a figure environment, an image-inclusion command or drawing code, so no image is delivered and the package declares no asset dependency.  Editorial formatting only, in addition: an AMS \texttt{amsart} preamble that restates the article's own declarations and macros used by the retained lines, namely the environments \texttt{lemma} on separate counters, together with the standard packages those lines need.  The retained environments print in this document as Lemma 1 in order of appearance, whereas the article prints the same items with the numbering of its own full document.  The complete native source of the article as distributed in the arXiv e-print for this exact version is bundled unchanged as \texttt{sources/199838/source0.tex}.
\begin{lemma}\label{uniformizador de L}
     Let $f\in K$ such that $v(f)=-s$ for some $s\in \mathbb{N}$ prime to $p$. Consider the Artin--Schreier extension $L=K_{\wp, f}$. Let $x\in L$ be such that $x^{p}-x=f$ and $a,b\in \mathbb{Z}$ be such that $-sa+pb=1$.
    Then for any uniformizer $t$ of $K$, $\pi=x^{a}t^{b}\in L$ is a uniformizer of $L$.
\end{lemma}

\begin{lemma}\label{t em funcao de pi}
    Let $t$ be a uniformizer of $K$ and let $s\in \mathbb{N}$ be prime to $p$. Let $L\defeq K_{\wp, t^{-s}}$ with $v_{L}$ the normalized valuation of $L$. Let $x\in L$ be such that $x^{p}-x=f$ and $\pi$ be the uniformizer obtained in Lemma \ref{uniformizador de L} for $f=t^{-s}$. Then for every integer $m>0$ prime to $p$,
    \[
    \frac{1}{t^{m}}=\frac{1}{\pi^{mp}}+g\in L
    \]
    where $g\in L$ is such that $v_{L}(g)=(p-1)s-mp$.
\end{lemma}

\begin{proof}
    This is equivalent to proving that $v_{L}\left(\frac{1}{t^{m}}-\frac{1}{\pi^{mp}}\right)=(p-1)s-mp$.
    Let us start by calculating $v_{L}(t^{m}-\pi^{mp})$. Observe that
    \begin{align*}
        t^{m}-\pi^{mp} &=t^{m}-(x^{p})^{mu}t^{mpv}=t^{m}-\left(\frac{1}{t^{s}}+x\right)^{mu}t^{mpv}\;
        \\
        &=t^{m}-t^{mpv}\sum_{n\geq 0}\binom{mu}{n}\frac{x^{n}}{t^{(msu-sn)}}
        \\
        &=t^{m}-t^{mpv}\frac{1}{t^{smu}}-t^{mpv}\sum_{n\geq 1}\binom{mu}{n}\frac{x^{n}}{t^{(msu-sn)}}
        \\
        &=t^{m}-t^{m(-su+pv)}-t^{mpv}\sum_{n\geq 1}\binom{mu}{n}\frac{x^{n}}{t^{(msu-sn)}}
        \\
        &=-t^{mpv}\sum_{n\geq 1}\binom{mu}{n}\frac{x^{n}}{t^{(msu-sn)}}\;.
    \end{align*}
    To calculate $v_{L}(\sum_{n\geq 1}\binom{mu}{n}\frac{x^{n}}{t^{(msu-sn)}})$, observe that for all $\binom{mu}{n}$ not divisible by $p$, we have
    \[
    v_{L}\left(\binom{mu}{n}\frac{x^{n}}{t^{(msu-sn)}}\right)=-pmsu+sn(p-1)\;.
    \]
    Since all of these valuations are different as $n$ varies, we have that
    \begin{align*}
        v_{L}\left(\sum_{n\geq 1}\binom{mu}{n}\frac{x^{n}}{t^{(msu-sn)}}\right) &=\text{min}\{-pmsu+sn(p-1)\}
        \\
        &=-pmsu+s(p-1)\;,
    \end{align*}
    where the minimun is taken among $1\leq n\leq mu$ such that $p\nmid \binom{mu}{n}$, and achieved at $n=1$. Therefore, 
    \begin{align*}
        v_{L}(t^{m}-\pi^{mp}) &=v_{L}(-t^{mpv}\sum_{n\geq 1}\binom{mu}{n}\frac{x^{n}}{t^{(msu-sn)}})
        \\
        &=mp^{2}v-pmsu+s(p-1)
        \\
        &=mp(-su+pv)+s(p-1)
        \\
        &=mp+s(p-1)\;.
    \end{align*}
    We then obtain that
    \begin{align*}
        v_{L}\left(\frac{1}{t^{m}}-\frac{1}{\pi^{mp}}\right)&=v_{L}\left(\frac{\pi^{mp}-t^{m}}{t^{m}\pi^{mp}}\right)
        \\
        &=v_{L}(\pi^{mp}-t^{m})-v_{L}(t^{m}\pi^{mp})
        \\
        &=mp+(p-1)s-2mp
        \\
        &=(p-1)s-mp,
    \end{align*}
    finishing the proof.
\end{proof}

\end{document}
